% SOURCE_STATEMENT: "For multipartite quantum systems with more than two parties, can the bipartite data-hiding results for separable states be generalized, including asymptotically optimal information efficiency and the best hiding-quality scaling?"
% Preparation and final self-audit: CONTRIBUTING.md prompts C and D; source/audit artifacts retained in the local physics_statements_2026_10_11 batch.
% !TeX program = lualatex
\documentclass[11pt]{article}
\usepackage[margin=25mm]{geometry}
\usepackage{fontspec}
\setmainfont{Latin Modern Roman}
\usepackage{amsmath,amssymb,mathtools}
\usepackage{unicode-math}
\setmathfont{Latin Modern Math}
\usepackage{xurl,hyperref}
\hypersetup{colorlinks=true,urlcolor=blue}
\setcounter{secnumdepth}{0}
\setlength{\parindent}{0pt}
\setlength{\parskip}{5pt}
\setlength{\emergencystretch}{3em}
\newcommand{\sref}[2]{\hyperlink{source:#1:#2}{[S#2]}}
\newcommand{\cataloguescope}[1]{\paragraph{Recorded target.}#1}
\newcommand{\R}{\mathbb R}
\newcommand{\C}{\mathbb C}
\newcommand{\N}{\mathbb N}
\newcommand{\Z}{\mathbb Z}
\newcommand{\Q}{\mathbb Q}
\newcommand{\F}{\mathbb F}
\newcommand{\eps}{\varepsilon}
\begin{document}
% UnsolvedMath ID 30002674; source catalogue code OWR-13111-004
\setcounter{section}{30002673}
\section{Multipartite Quantum Data Hiding with Separable States}\label{30002674}
{\small UnsolvedMath ID 30002674 · Mathematical Physics}
\subsection{Definitions and mathematical statement}

\paragraph{Shared notation.}$\mathbb N=\{1,2,\ldots\}$; the standard number fields have their usual meanings. All additional conventions are stated locally.


\paragraph{Definitions and conventions.}
A multipartite state is a positive trace-one operator on $\mathcal H_1\otimes\cdots\otimes\mathcal H_m$, with $m>2$. A fully separable state is a convex combination of product states. Data hiding compares distinguishability by unrestricted measurements with distinguishability under an allowed local or separable measurement class. Information efficiency measures how much classical information is encoded relative to quantum-system size; hiding quality measures the gap between global and restricted distinguishability.

\paragraph{Target retained from the canonical source catalogue.}
\paragraph{Catalogue wording retained for traceability.}
For multipartite quantum systems with more than two parties, can the bipartite data-hiding results for separable states be generalized, including asymptotically optimal information efficiency and the best hiding-quality scaling?

\paragraph{Source scope and qualification.}
The restriction on allowed measurements and the asymptotic family are essential and must be read from the source bipartite results. LOCC and all separable measurements are different measurement classes; they are not identified here.
\subsection{Short English statement}
Can separable quantum states hide information efficiently among more than two parties, with optimal hiding strength?
\subsection{Sources}
\begin{itemize}
\item[\hypertarget{source:30002674:1}{S1}] ULAM.AI and the UnsolvedMath Contributors, UnsolvedMath, record 30002674 (OWR-13111-004). Catalogue curation under CC BY 4.0; definitions and English exposition added locally.
\url{https://huggingface.co/datasets/ulamai/UnsolvedMath}.
\item[\hypertarget{source:30002674:2}{S2}] New Horizons in Statistical Decision Theory (2014), Oberwolfach Reports, DOI 10.4171/owr/2014/41, pp. 2297--2352. Primary report accessed 2026-10-11; relevant PDF pages 50.
\url{https://doi.org/10.4171/owr/2014/41}.
\item[\hypertarget{source:30002674:3}{S3}] Hayden--Leung--Smith, multiparty data hiding (catalogue-listed context; not independently consulted in this preparation).
\url{https://doi.org/10.1103/PhysRevA.71.062339}.
\end{itemize}
\label{end:30002674}
\end{document}
